\documentclass{article}
% Source: Topic_03_Teacher_Edition.pdf, PDF pages 101-102 (printed pages 177A-177B): Topic Assessment Form A with the TE's marked answers, Form B with marked answers, Item Analysis for both forms. Cross-checked against the Realize download aga_24_a2_na_03_ttb.pdf (Form B text layer).
\usepackage{amsmath}
\usepackage{amssymb}
\newenvironment{assessment-meta}{}{}   % topic, title, forms, item count, printed pages, sources
\newenvironment{item-analysis}{}{}     % Savvas's Item / DOK / Skills Review and Practice table (same for both forms)
\newenvironment{ta-item}[2]{}{}        % #1 = item number, #2 = form (A or B)
\newcommand{\answer}[1]{}              % the TE's marked answer, inside a ta-item
\newcommand{\placeholder}[2]{}         % #1 = type (graph, table, figure, number-line), #2 = complete description
\newcommand{\uncertain}[2]{#1}         % #1 = best reading, #2 = alternatives and why
\newcommand{\te}[1]{}                  % transcriber's note (a discrepancy, a publisher misprint, a computed check)
\begin{document}
\begin{assessment-meta}
topic: 3
title: Polynomial Functions
forms: A and B (parallel item for item)
items: 14
printed: 177A-177B
sources: Topic_03_Teacher_Edition.pdf PDF pages 101-102 (TE page heads "TOPIC 3 Polynomial Functions", "TOPIC ASSESSMENT Form A" / "Form B"; footers "TOPIC 3 | 177A | TOPIC ASSESSMENT" and "TOPIC 3 | 177B | TOPIC ASSESSMENT"); student pages headed "3 Topic Assessment Form A" / "3 Topic Assessment Form B", each 2 pages ("enVision Algebra 2 1of2 / 2of2 Assessment Sourcebook"); Realize download aga_24_a2_na_03_ttb.pdf (Form B, text only) for cross-checking.
\end{assessment-meta}
\te{TE note under Form A (177A): "Topic Assessment Assess students' understanding of topic concepts and skills using the Topic Assessment found at SavvasRealize.com. These auto-scored online assessments provide students with a breadth of technology-enhanced item types. There are two versions of the Topic Assessment, Form A and Form B. These two versions, available in print and at SavvasRealize.com, are parallel tests that assess the same content item for item. The Item Analysis chart on the next page can be used for both versions."}
% Source page 102: 177B Item Analysis
\begin{item-analysis}
Item & DOK & Skills Review and Practice \\
1 & 1 & A74, A75 \\
2 & 2 & F62 \\
3 & 2 & A20 \\
4 & 1 & F62 \\
5 & 2 & A74, A75 \\
6 & 2 & F61 \\
7 & 3 & G100 \\
8 & 2 & F61 \\
9 & 2 & A73 \\
10 & 2 & A75 \\
11 & 2 & A73 \\
12 & 1 & F63 \\
13 & 2 & F94 \\
14 & 2 & F73 \\
\end{item-analysis}
% Source page 101: 177A Topic Assessment Form A
\begin{ta-item}{1}{A}
What is the remainder when (f(x)=2x^4+x^3-8x-1) is divided by (x-2)?
A. (-23)
B. (23)
C. (-3)
D. (3)
\answer{1. B (23)}
\te{computed: f(2)=32+8-16-1=23; agrees. Choices print in two columns (A, B left; C, D right).}
\end{ta-item}
\begin{ta-item}{2}{A}
Describe the coefficients of a cubic polynomial function (f) that has one rational zero, (a), and two irrational zeros, (\sqrt{b}) and (-\sqrt{b}), where (b) is rational.
A. Some are irrational numbers.
B. They are all complex numbers.
C. They are all rational numbers.
D. They are all irrational numbers.
\answer{2. C. They are all rational numbers.}
\te{computed: (x-a)(x^2-b)=x^3-ax^2-bx+ab, all rational; agrees for a monic cubic (or any rational leading coefficient), which the item intends; with an irrational leading coefficient the coefficients need not be rational (Codex review 2026-10-31).}
\end{ta-item}
\begin{ta-item}{3}{A}
Use the Binomial Theorem to expand ((x-2)^6).
A. (x^6-2x^5+4x^4-8x^3+16x^2-32x+64)
B. (x^6-12x^5+24x^4-36x^3+48x^2-60x+12)
C. (x^6-12x^5+60x^4-160x^3+240x^2-192x+64)
D. (x^6-32x^5+16x^4-8x^3+4x^2-2x+64)
\answer{3. C. (x^6-12x^5+60x^4-160x^3+240x^2-192x+64)}
\te{computed: agrees.}
\end{ta-item}
\begin{ta-item}{4}{A}
Simplify ((x^2+2x)(x^2+x+5)). Write the simplified polynomial in standard form.
\placeholder{figure}{Drag-and-drop tiles, two rows: (x^4), (2x^3), (3x^3), (5x^2), (7x^2) / (10x), (10). Answer row of four boxes: [ ] + [ ] + [ ] + [ ].}
\answer{4. (x^4) + (3x^3) + (7x^2) + (10x)}
\te{computed: x^4+3x^3+7x^2+10x; agrees.}
\end{ta-item}
\begin{ta-item}{5}{A}
If (x-3) is a factor of (P(x)=x^3+x^2-17x+15), select the two factors that express (P(x)) as a product.
A. ((x-3))
B. ((x^2-4x-5))
C. ((x^2+4x+5))
D. ((x^2+4x-5))
E. not a factor
\answer{5. A, D}
\te{computed: (x-3)(x^2+4x-5)=x^3+x^2-17x+15; agrees.}
\end{ta-item}
\begin{ta-item}{6}{A}
Sketch a graph of the polynomial function (f(x)=x^3-2x^2).
\placeholder{graph}{Answer graph: gray grid, -4 to 4 on each axis, one unit per grid square; the axes (labeled x and y, arrowheads at both ends), the tick labels (-2 and 2 on the x-axis, 2 and -2 on the y-axis, origin O) and the curve are printed in magenta; the curve is the TE's answer. Drawn curve: comes up from an arrow at the bottom near (-0.6,-4), crosses the x-axis at about the origin, rises to a local maximum near (0.4,0.6), falls through the x-axis near x=1.2 to a local minimum near (2,-2), then rises through the x-axis near x=2.9 to an arrow at the top near (3.3,4).}
\answer{6. The magenta curve described above.}
\te{computed: f(x)=x^2(x-2) touches the x-axis at 0 (local maximum (0,0)), has a local minimum at (4/3,-32/27)\approx(1.33,-1.19), and crosses at x=2. The printed sketch does not match: it crosses near 0 with a maximum above the axis and has three x-intercepts (about 0, 1.2, 2.9) and a minimum near (2,-2). Coordinates read from the 600 dpi crop are approximate.}
\end{ta-item}
\begin{ta-item}{7}{A}
Lucy cuts 4 squares with side length (x) in. from the corners of a 6-by-9-in. piece of paper. She folds the paper to make a tray that is (x) in. high. Complete the function for the tray's volume (V) in terms of (x).
(V(x)=) [ ] (x^3-) [ ] (x^2+) [ ] (x)
\answer{7. 4, 30, 54 (V(x)=4x^3-30x^2+54x)}
\te{computed: x(6-2x)(9-2x)=4x^3-30x^2+54x; agrees.}
\end{ta-item}
\begin{ta-item}{8}{A}
Use polynomial identities to factor (64+27a^3). Select all factors.
A. ((4-3a))
B. ((4+3a))
C. ((16+12a+9a^2))
D. ((16-12a+9a^2))
E. ((16-12a+3a^2))
\answer{8. B, D}
\te{computed: (4+3a)(16-12a+9a^2); agrees.}
\end{ta-item}
\begin{ta-item}{9}{A}
The volume of a cube with length (x) is (V(x)=x^3). The volume of a sphere with radius (x) is shown in the graph.
\placeholder{graph}{Prompt graph (black): first quadrant only; x-axis labeled x with labels 0 and 1, grid spacing 0.25 (window 0 to 2); y-axis labeled y with labels 0, 5, 10, 15, grid spacing 2.5 (window 0 to 20). One increasing curve starting at (0,0), flat near 0, passing about (1, 4.2), ending in an arrow near (1.65, 19); consistent with (V=\frac{4}{3}\pi x^3).}
Over the interval ([0.5, 1]), which volume is increasing faster?
The volume of the [sphere / cube] is increasing faster.
\answer{9. sphere}
\te{computed: average rate on [0.5,1]: cube (1-0.125)/0.5=1.75; sphere (\frac{4}{3}\pi(0.875)/0.5\approx7.33); agrees.}
\end{ta-item}
\begin{ta-item}{10}{A}
Use synthetic division to divide (2x^4-3x^3+2x^2-8x-1) by (x-1). Complete the quotient.
\placeholder{figure}{Drag-and-drop tiles, two rows: (2x^4), (2x^3), (-x^2), (-6x^2), (x) / (8x), (-7), (-16), (-\frac{8}{x-1}), (-\frac{15}{x-1}). Answer row: [ ] + [ ] + [ ] + [ ] + [ ].}
\answer{10. (2x^3) + (-x^2) + (x) + (-7) + (-\frac{8}{x-1})}
\te{computed: 2, -1, 1, -7, remainder -8; agrees.}
\end{ta-item}
\begin{ta-item}{11}{A}
Describe the behavior at each zero of (f(x)=(x-5)(x-1)(x+4)).
A. The graph touches the (x)-axis at 4 and crosses the (x)-axis at (-1) and (-5).
B. The graph touches the (x)-axis at 4, (-5), and (-1).
C. The graph crosses the (x)-axis at 5, 1, and (-4).
D. The graph touches the (x)-axis at (-4) and crosses the (x)-axis at 1 and 5.
\answer{11. C. The graph crosses the (x)-axis at 5, 1, and (-4).}
\te{computed: three simple zeros; agrees.}
\end{ta-item}
\begin{ta-item}{12}{A}
Identify the zeros of the function (f(x)=x(x-3)(x+1)).
[ ], [ ], [ ]
\answer{12. 0, 3, -1}
\end{ta-item}
\begin{ta-item}{13}{A}
Use the binomial theorem to expand ((2a+2b)^5). Select all of the terms below that are in the standard form of the expansion. (Not all terms are given.)
A. (5a^4b)
B. (384a^3b^3)
C. (160ab^4)
D. (16a^2b^2)
E. (32a^5)
F. (20a^2b^3)
\answer{13. C, E}
\te{computed: 32a^5+160a^4b+320a^3b^2+320a^2b^3+160ab^4+32b^5; agrees.}
\end{ta-item}
\begin{ta-item}{14}{A}
How does the graph of the function (f(x)=2x^4-3) differ from the graph of its parent function? Select all the transformations.
A. The graphs are the same.
B. Subtracting 3 translates the graph down 3 units.
C. The leading coefficient, 2, stretches the graph vertically.
D. Subtracting 3 translates the graph up 3 units.
E. Subtracting 3 translates the graph right 3 units.
\answer{14. B, C}
\end{ta-item}
% Source page 102: 177B Topic Assessment Form B
\begin{ta-item}{1}{B}
What is the remainder when (f(x)=x^4-8x^2+10x+22) is divided by (x+3)?
A. (-107)
B. (-17)
C. (61)
D. (1)
\answer{1. D (1)}
\te{computed: f(-3)=81-72-30+22=1; agrees. Choices print in two columns (A, B left; C, D right).}
\end{ta-item}
\begin{ta-item}{2}{B}
Describe the coefficients of a cubic polynomial function (f) that has one rational zero, (a), and two complex numbers (i\sqrt{b}) and (-i\sqrt{b}), where (b) is rational.
A. None are real numbers.
B. They are all real numbers.
C. Only one is a real number.
D. Only two are real numbers.
\answer{2. B. They are all real numbers.}
\te{computed: (x-a)(x^2+b)=x^3-ax^2+bx-ab, all rational (so real); agrees. The Realize text layer drops the radicals ("ib and - i b"); the TE image shows i\sqrt{b}.}
\end{ta-item}
\begin{ta-item}{3}{B}
Use the binomial theorem to expand ((2x-2)^6).
A. (2x^6-12x^5+30x^4-20x^3+30x^2-12x+2)
B. (2x^6-24x^5+120x^4-320x^3+480x^2-384x+128)
C. (64x^6-192x^5+240x^4-160x^3+60x^2-12x+1)
D. (64x^6-384x^5+960x^4-1280x^3+960x^2-384x+64)
\answer{3. D. (64x^6-384x^5+960x^4-1280x^3+960x^2-384x+64)}
\te{computed: 64(x-1)^6; agrees.}
\end{ta-item}
\begin{ta-item}{4}{B}
Simplify ((m+3)(m^2-2m+2)). Write the simplified polynomial in standard form.
\placeholder{figure}{Drag-and-drop tiles, two rows: (5), (6), (m^3), (m^2), (-m^2) / (2m), (-4m), (10m). Answer row of four boxes: [ ] + [ ] + [ ] + [ ].}
\answer{4. (m^3) + (m^2) + (-4m) + (6)}
\te{computed: m^3+m^2-4m+6; agrees.}
\end{ta-item}
\begin{ta-item}{5}{B}
If (x+5) is a factor of (P(x)=x^3+4x^2+2x+35), select the two factors that express (P(x)) as a product.
A. ((x+5))
B. ((x^2+x+7))
C. ((x^2+x-3))
D. ((x^2-x+7))
E. not a factor
\answer{5. A, D}
\te{computed: (x+5)(x^2-x+7)=x^3+4x^2+2x+35; agrees.}
\end{ta-item}
\begin{ta-item}{6}{B}
Sketch a graph of the polynomial function (f(x)=x^3+\frac{1}{2}x^2-x-\frac{1}{2}).
\placeholder{graph}{Answer graph: black grid about -4 to 4 on each axis, one unit per grid square; black axes labeled x and y with arrowheads at both ends; tick labels -2 and 2 on the x-axis, 2 and -2 on the y-axis, origin O. Magenta curve (the printed answer): rises from an arrow at the bottom near (-1.9,-4), reaches the x-axis near x=-1 with a small local maximum just above the axis near (-0.7,0.1), crosses back down near x=-0.4, passes about (0,-0.5) to a local minimum near (0.45,-0.75), crosses the x-axis at x=1, and rises to an arrow at the top near (1.7,4).}
\answer{6. The magenta curve described above.}
\te{computed: f(x)=(x+1)(x+\frac{1}{2})(x-1), zeros -1, -1/2, 1, local maximum about (-0.77,0.11), local minimum about (0.43,-0.76), y-intercept -1/2; the sketch is consistent. Coordinates read from the 600 dpi crop are approximate.}
\end{ta-item}
\begin{ta-item}{7}{B}
Teo cuts 4 squares with side length (x) in. from the corners of a 10-by-15-in. piece of paper. He folds the paper to make a tray that is (x) in. high. Complete the function for the tray's volume (V).
(V(x)=) [ ] (x^3-) [ ] (x^2+) [ ] (x)
\answer{7. 4, 50, 150 (V(x)=4x^3-50x^2+150x)}
\te{computed: x(10-2x)(15-2x)=4x^3-50x^2+150x; agrees.}
\end{ta-item}
\begin{ta-item}{8}{B}
Use polynomial identities to factor (64x^3-1). Select all factors.
A. ((4x+1))
B. ((4x-1))
C. ((16x^2-4x+1))
D. ((16x^2+4x+1))
E. ((4x^2+4x+1))
\answer{8. B, D}
\te{computed: (4x-1)(16x^2+4x+1); agrees.}
\end{ta-item}
\begin{ta-item}{9}{B}
The volume of a cube with length (x) is (V(x)=x^3). The volume of a sphere with radius (\frac{1}{2}x) is shown in the graph.
\placeholder{graph}{Prompt graph (black), the same graph as Form A item 9: first quadrant only; x-axis labeled x with labels 0 and 1, grid spacing 0.25 (window 0 to 2); y-axis labeled y with labels 0, 5, 10, 15, grid spacing 2.5 (window 0 to 20). One increasing curve starting at (0,0), flat near 0, passing about (1, 4), ending in an arrow near (1.65, 19).}
Over the interval ([1, 2]), which volume is increasing faster?
The volume of the [sphere / cube] is increasing faster.
\answer{9. cube}
\te{computed: a sphere of radius x/2 has V=\frac{\pi}{6}x^3\approx0.52x^3, so on [1,2] the cube's average rate is 7 and the sphere's about 3.67; the printed key (cube) agrees with the text. The printed graph, however, is the Form A curve (about 4.2 at x=1, i.e. \frac{4}{3}\pi x^3, radius x), not \frac{\pi}{6}x^3 (0.52 at x=1); read from the graph, the sphere would be faster. The graph also ends near x=1.65, short of the interval's right end 2. Publisher misprint in the figure.}
\end{ta-item}
\begin{ta-item}{10}{B}
Use synthetic division to divide (x^4-5x^3+2x^2+6x+6) by (x-2). Complete the quotient.
\placeholder{figure}{Drag-and-drop tiles, two rows: (x^4), (x^3), (-3x^2), (-7x^2), (-4x) / (16x), (-2), (-26), (\frac{2}{x-2}), (\frac{58}{x-2}). Answer row: [ ] + [ ] + [ ] + [ ] + [ ].}
\answer{10. (x^3) + (-3x^2) + (-4x) + (-2) + (\frac{2}{x-2})}
\te{computed: 1, -3, -4, -2, remainder 2; agrees.}
\end{ta-item}
\begin{ta-item}{11}{B}
Describe the behavior at each zero of the function (f(x)=(x-5)(x-1)^2).
A. The graph touches the (x)-axis at 1 and crosses the (x)-axis at 5.
B. The graph touches the (x)-axis at (-5) and 1.
C. The graph touches the (x)-axis at (-1) and crosses the (x)-axis at (-5).
D. The graph crosses the (x)-axis at (-1) and 11.
\answer{11. A. The graph touches the (x)-axis at 1 and crosses the (x)-axis at 5.}
\te{computed: double zero at 1, simple zero at 5; agrees.}
\end{ta-item}
\begin{ta-item}{12}{B}
Identify the zeros of the function (f(x)=x(x-1)(x+2)).
[ ], [ ], [ ]
\answer{12. 0, 1, -2}
\end{ta-item}
\begin{ta-item}{13}{B}
Use the binomial theorem to expand ((3a+3b)^4). Select all of the terms below that are in the standard form of the expansion. (Not all terms are given.)
A. (486a^2)
B. (486a^2b^2)
C. (81ab^3)
D. (324a^3b)
E. (324ab^3)
F. (729a^2b^3)
\answer{13. B, D, E}
\te{computed: 81a^4+324a^3b+486a^2b^2+324ab^3+81b^4; agrees.}
\end{ta-item}
\begin{ta-item}{14}{B}
How does the graph of the function (f(x)=2x^3-1) differ from the graph of its parent function? Select all the transformations.
A. The graphs are the same.
B. Subtracting 1 translates the graph right 3 units.
C. The leading coefficient, 2, stretches the graph vertically.
D. Subtracting 3 translates the graph left 3 units.
E. Subtracting 1 translates the graph down 1 unit.
\answer{14. C, E}
\end{ta-item}
\end{document}
